\section{Рабочий проект}
\subsection{Описание модулей разработанной программной системы}

Разработанная программно-информационная система регистрации данных пользователей реализована в виде одностраничного web-приложения (Single Page Application), построенного с использованием фреймворка Vue~3 и языка TypeScript. Архитектура приложения основана на доменном подходе, обеспечивающем разделение функциональности по предметным областям и снижение связанности между отдельными модулями системы.

Каждый домен включает компоненты пользовательского интерфейса, бизнес-логику, средства доступа к данным и описание используемых типов данных. Такой подход упрощает сопровождение и расширение системы, а также повышает удобство тестирования отдельных функциональных частей приложения.

Основные домены разработанной системы представлены в таблице \ref{domains:table}.

\newpage

\renewcommand{\arraystretch}{1.0}

\begin{xltabular}[ht]
	{\textwidth}{|p{2.5cm}|>{\RaggedRight\arraybackslash}p{4cm}|>{\RaggedRight\arraybackslash}X|}
	\caption{Описание доменов программной системы\label{domains:table}}\\
	\hline
	\centrow Домен &
	\centrow Назначение &
	\centrow Основные функции\\
	\hline
	\endfirsthead
	
	Auth &
	Аутентификация и авторизация &
	Вход пользователей в систему, завершение сеанса работы, получение информации о текущем пользователе, определение роли пользователя\\
	\hline
	
	Customers &
	Управление клиентами &
	Просмотр, поиск, создание, редактирование и удаление карточек клиентов\\
	\hline
	
	Vehicles &
	Учёт автомобилей &
	Хранение сведений об автомобилях клиентов, регистрация новых автомобилей, редактирование данных\\
	\hline
	
	Requests &
	Управление заявками &
	Создание и обработка заявок, изменение статусов, назначение мастеров, контроль жизненного цикла заявок\\
	\hline
	
	Users &
	Управление пользователями &
	Создание учётных записей, назначение ролей, редактирование профилей пользователей\\
	\hline
	
	Audit &
	Журнал изменений &
	Просмотр истории изменений, фильтрация событий аудита, анализ действий пользователей\\
	\hline
	
\end{xltabular}

\renewcommand{\arraystretch}{1.0}

Архитектурная организация системы обеспечивает независимость доменов друг от друга. Взаимодействие между слоями осуществляется только в соответствии с установленными правилами зависимостей, что позволяет минимизировать влияние изменений в одном модуле на остальные части приложения.

На рисунке \ref{project-structure:image} представлена структура каталогов разработанной программной системы.

\begin{figure}[H]
	\centering
	\includegraphics[width=1.0\linewidth]{project-structure}
	\caption{Структура каталогов программной системы}
	\label{project-structure:image}
\end{figure}

Представленная организация исходного кода обеспечивает локализацию бизнес-логики внутри соответствующих доменов и позволяет независимо развивать отдельные функциональные подсистемы приложения.

\newpage
\subsection{Описание компонентов и классов приложения}

Внутри каждого домена используется единая структура организации исходного кода. Основные компоненты представлены в таблице \ref{components:table}.

\renewcommand{\arraystretch}{0.85}

\begin{xltabular}{\textwidth}{|p{2.5cm}|p{3cm}|X|X|}
	\caption{Описание компонентов домена\label{components:table}}\\
	\hline
	\centrow Компонент &
	\centrow Назначение &
	\centrow Описание &
	\centrow Пример использования\\
	\hline
	\endfirsthead
	
	\caption*{Продолжение таблицы \ref{components:table}}\\
	\hline
	\centrow Компонент &
	\centrow Назначение &
	\centrow Описание &
	\centrow Пример использования\\
	\hline
	\finishhead
	
	api &
	Доступ к данным &
	Выполнение запросов к Supabase &
	requests-api.ts\\
	\hline
	
	store &
	Хранение состояния &
	Содержит доменное состояние и действия Pinia &
	requests/store.ts\\
	\hline
	
	composables &
	Бизнес-логика &
	Реализация переиспользуемых алгоритмов &
	useRequestWorkflow.ts\\
	\hline
	
	components &
	Компоненты интерфейса &
	Формы, таблицы, карточки и элементы управления &
	RequestForm.vue\\
	\hline
	
	views &
	Страницы приложения &
	Отображение законченных пользовательских сценариев &
	RequestsView.vue\\
	\hline
	
	types &
	Типизация &
	Описание интерфейсов и перечислений TypeScript &
	types.ts\\
	\hline
	
\end{xltabular}

Доступ к данным осуществляется исключительно через API-слой. Компоненты пользовательского интерфейса не взаимодействуют с Supabase напрямую. Это позволяет централизовать обработку ошибок и обеспечить единообразный механизм получения и изменения данных.

Направление зависимостей между слоями приложения представлено на рисунке \ref{layers:image}.

\begin{figure}[H]
	\centering
	\includegraphics[width=0.35\linewidth]{layers}
	\caption{Направление зависимостей между слоями приложения}
	\label{layers:image}
\end{figure}

Представленная схема отражает принцип инверсии зависимостей, согласно которому внутренние слои не зависят от внешних.

Пример реализации метода получения списка заявок приведён ниже.

\begin{lstlisting}[language=Java]
	export const requestsApi = {
		async getAll(): Promise<ServiceRequest[]> {
			const { data, error } = await supabase
			.from('service_requests')
			.select('*')
			
			if (error) {
				throw error
			}
			
			return data
		}
	}
\end{lstlisting}

Использование единой точки доступа к данным упрощает модификацию приложения при изменении используемых внешних сервисов и облегчает проведение модульного тестирования.

Схема взаимодействия компонентов системы при получении данных представлена на рисунке \ref{interaction:image}.

\begin{figure}[H]
	\centering
	\includegraphics[width=1.0\linewidth]{interaction}
	\caption{Взаимодействие компонентов приложения при получении данных}
	\label{interaction:image}
\end{figure}

Последовательность вызовов обеспечивает разделение ответственности между пользовательским интерфейсом, бизнес-логикой и средствами доступа к данным.

\subsection{Реализация ключевых алгоритмов}

Одним из ключевых алгоритмов системы является управление жизненным циклом заявок на обслуживание. Заявка может находиться только в одном из заранее определённых состояний, а переход между состояниями осуществляется в соответствии с установленными бизнес-правилами.

Взаимосвязи между основными сущностями предметной области представлены на рисунке \ref{er:image}.

\begin{figure}[H]
	\centering
	\includegraphics[width=0.7\linewidth]{er}
	\caption{Фрагмент схемы базы данных программной системы}
	\label{er:image}
\end{figure}

Каждый клиент может иметь несколько автомобилей, а для каждого автомобиля может быть зарегистрировано несколько заявок на обслуживание.

Жизненный цикл заявки реализован в виде конечного автомата. Допустимые переходы между статусами представлены на диаграмме \ref{workflow:image}.

\begin{figure}[H]
	\centering
	\includegraphics[width=0.6\linewidth]{workflow}
	\caption{Диаграмма жизненного цикла заявки}
	\label{workflow:image}
\end{figure}

Перечень допустимых переходов между состояниями представлен в таблице \ref{workflow:table}.\\

\begin{table}[H]
	\centering
	\caption{Допустимые переходы между статусами заявок}
	\label{workflow:table}
	
	\renewcommand{\arraystretch}{1.5} % увеличение высоты строк
	
	\begin{tabular}{|C{5cm}|C{8cm}|}
		\hline
		\textbf{Текущий статус} &
		\textbf{Допустимые переходы} \\
		\hline
		
		Черновик &
		Отправлена, Отменена \\
		\hline
		
		Отправлена &
		Назначена, Отменена \\
		\hline
		
		Назначена &
		В работе, Отменена \\
		\hline
		
		В работе &
		Завершена, Приостановлена \\
		\hline
		
		Приостановлена &
		В работе \\
		\hline
		
		Завершена &
		Переходы отсутствуют \\
		\hline
		
		Отменена &
		Переходы отсутствуют \\
		\hline
		
	\end{tabular}
\end{table}

Фрагмент программной реализации конечного автомата представлен ниже.

\begin{lstlisting}[language=Java]
	export const REQUEST_STATUS_TRANSITIONS = {
		DRAFT: ['SUBMITTED', 'CANCELLED'],
		SUBMITTED: ['ASSIGNED', 'CANCELLED'],
		ASSIGNED: ['IN_PROGRESS', 'CANCELLED'],
		IN_PROGRESS: ['COMPLETED', 'ON_HOLD'],
		ON_HOLD: ['IN_PROGRESS'],
		COMPLETED: [],
		CANCELLED: []
	}
\end{lstlisting}

Дополнительно выполняется проверка полномочий пользователя на выполнение операции изменения статуса. Проверка осуществляется с учётом роли пользователя, текущего состояния заявки и контекста её обработки.

Использование конечного автомата позволяет исключить возникновение некорректных состояний заявок и обеспечивает соблюдение регламентированных бизнес-процессов организации.

\subsection{Модульное тестирование}

Поскольку разработка программной системы осуществлялась без применения специализированных средств автоматизированного модульного тестирования, проверка корректности работы отдельных функциональных модулей выполнялась вручную посредством последовательного выполнения тестовых сценариев.

Основной целью модульного тестирования являлась проверка корректности работы отдельных компонентов системы независимо друг от друга.

Результаты проверки представлены в таблице \ref{module:table}.

\renewcommand{\arraystretch}{1.2}

\begin{xltabular}{\textwidth}
	{|p{3.5cm}|X|p{2.5cm}|}
	\caption{Результаты модульного тестирования\label{module:table}}\\
	\hline
	\centrow Проверяемый модуль &
	\centrow Проверяемая функциональность &
	\centrow Результат\\
	\hline
	
	Аутентификация &
	Проверка входа пользователя с корректными учётными данными &
	Успешно\\
	\hline
	
	Аутентификация &
	Проверка отказа во входе при вводе неверного пароля &
	Успешно\\
	\hline
	
	Клиенты &
	Создание карточки клиента &
	Успешно\\
	\hline
	
	Клиенты &
	Редактирование данных клиента &
	Успешно\\
	\hline
	
	Автомобили &
	Регистрация нового автомобиля &
	Успешно\\
	\hline
	
	Автомобили &
	Изменение сведений об автомобиле &
	Успешно\\
	\hline
	
	Заявки &
	Создание заявки на обслуживание &
	Успешно\\
	\hline
	
	Заявки &
	Изменение статуса заявки &
	Успешно\\
	\hline
	
	Пользователи &
	Назначение роли пользователю &
	Успешно\\
	\hline
	
	Журнал аудита &
	Отображение истории изменений &
	Успешно\\
	\hline
	
\end{xltabular}

\renewcommand{\arraystretch}{1.0}

\subsection{Системное тестирование}

Системное тестирование проводилось путём проверки основных пользовательских сценариев работы программно-информационной системы через веб-интерфейс приложения. Целью тестирования являлось подтверждение корректности функционирования реализованных бизнес-процессов и соответствия системы требованиям технического задания.

\subsubsection{Тестовый сценарий: Авторизация пользователя}

\textbf{Предусловие:} пользователь зарегистрирован в системе и обладает действующей учётной записью.

\textbf{Тестовый случай:} пользователь вводит адрес электронной почты и пароль, после чего нажимает кнопку "<Войти">.

\textbf{Ожидаемый результат:} система выполняет проверку введённых данных посредством обращения к Supabase, формирует пользовательскую сессию и перенаправляет пользователя на главную панель приложения в соответствии с назначенной ролью.

\textbf{Результат:} тест пройден успешно.

\begin{figure}[H]
	\centering
	\includegraphics[width=0.85\linewidth]{login}
	\caption{Экран авторизации пользователя}
	\label{login:image}
\end{figure}


\subsubsection{Тестовый сценарий: Отображение главной панели системы}

\textbf{Предусловие:} пользователь успешно прошёл процедуру авторизации.

\textbf{Тестовый случай:} после входа в систему осуществляется переход на главную панель приложения.

\textbf{Ожидаемый результат:} пользователю отображается главная панель с основными разделами системы и сводной информацией, доступной в соответствии с его ролью.

\textbf{Результат:} тест пройден успешно.

\begin{figure}[H]
	\centering
	\includegraphics[width=0.95\linewidth]{dashboard}
	\caption{Главная панель программной системы}
	\label{dashboard:image}
\end{figure}

\subsubsection{Тестовый сценарий: Управление клиентами}

\textbf{Предусловие:} пользователь обладает правами доступа к разделу клиентов.

\textbf{Тестовый случай:} пользователь открывает раздел "<Клиенты">, осуществляет просмотр существующих записей и выполняет добавление либо редактирование данных клиента.

\textbf{Ожидаемый результат:} изменения корректно сохраняются и отображаются в интерфейсе приложения.

\textbf{Результат:} тест пройден успешно.

\begin{figure}[H]
	\centering
	\includegraphics[width=0.95\linewidth]{customers}
	\caption{Интерфейс управления клиентами}
	\label{customers:image}
\end{figure}

\subsubsection{Тестовый сценарий: Добавление клиента}

\textbf{Предусловие:} пользователь авторизован.

\textbf{Тестовый случай:} открыть раздел клиентов и нажать "<Добавить клиента">.

\textbf{Ожидаемый результат:} новый клиент сохраняется и отображается в списке.

\textbf{Результат:} тест пройден успешно.

\begin{figure}[H]
	\centering
	\includegraphics[width=0.8\linewidth]{customer-create}
	\caption{Добавление клиента}
	\label{customer-create:image}
\end{figure}

\subsubsection{Тестовый сценарий: Редактирование данных клиента}

\textbf{Предусловие:} в системе существует клиент.

\textbf{Тестовый случай:} изменить контактный телефон клиента.

\textbf{Ожидаемый результат:} изменения сохраняются и отображаются после обновления страницы.

\textbf{Результат:} тест пройден успешно.

\begin{figure}[H]
	\centering
	\includegraphics[width=0.8\linewidth]{customer-edit}
	\caption{Редактирование данных клиента}
	\label{customer-edit:image}
\end{figure}

\subsubsection{Тестовый сценарий: Управление автомобилями}

\textbf{Предусловие:} в системе зарегистрированы клиенты.

\textbf{Тестовый случай:} пользователь открывает раздел автомобилей и выполняет добавление либо редактирование сведений об автомобиле клиента.

\textbf{Ожидаемый результат:} информация об автомобиле сохраняется в базе данных и становится доступной для дальнейшего использования при оформлении заявок.

\textbf{Результат:} тест пройден успешно.

\begin{figure}[H]
	\centering
	\includegraphics[width=0.95\linewidth]{vehicles}
	\caption{Интерфейс управления автомобилями}
	\label{vehicles:image}
\end{figure}

\subsubsection{Тестовый сценарий: Добавление автомобиля}

\textbf{Предусловие:} в системе существует клиент.

\textbf{Тестовый случай:} создать новую запись автомобиля.

\textbf{Ожидаемый результат:} автомобиль появляется в списке и становится доступным при оформлении заявок.

\textbf{Результат:} тест пройден успешно.

\begin{figure}[H]
	\centering
	\includegraphics[width=0.8\linewidth]{vehicle-create}
	\caption{Добавление автомобиля}
	\label{vehicle-create:image}
\end{figure}

\subsubsection{Тестовый сценарий: Просмотр списка заявок}

\textbf{Предусловие:} в базе данных присутствуют зарегистрированные заявки.

\textbf{Тестовый случай:} пользователь переходит в раздел «Заявки».

\textbf{Ожидаемый результат:} система отображает список заявок с указанием основных атрибутов: номера, клиента, автомобиля, текущего статуса и ответственного исполнителя.

\textbf{Результат:} тест пройден успешно.

\begin{figure}[H]
	\centering
	\includegraphics[width=0.95\linewidth]{requests}
	\caption{Просмотр списка заявок}
	\label{requests:image}
\end{figure}

\subsubsection{Тестовый сценарий: Создание заявки}

\textbf{Предусловие:} в системе существуют клиент и автомобиль.

\textbf{Тестовый случай:} заполнить форму заявки.

\textbf{Ожидаемый результат:} заявка сохраняется со статусом "<Черновик">.

\textbf{Результат:} тест пройден успешно.

\begin{figure}[H]
	\centering
	\includegraphics[width=0.8\linewidth]{request-create}
	\caption{Создание заявки}
	\label{request-create:image}
\end{figure}

\begin{figure}[H]
	\centering
	\includegraphics[width=0.95\linewidth]{request-form}
	\caption{Пример созданной заявки}
	\label{requestform:image}
\end{figure}

\subsubsection{Тестовый сценарий: Редактирование заявки}

\textbf{Предусловие:} в системе существуют заявка.

\textbf{Тестовый случай:} изменить описание заявки.

\textbf{Ожидаемый результат:} изменения сохраняются.

\textbf{Результат:} тест пройден успешно.

\begin{figure}[H]
	\centering
	\includegraphics[width=0.8\linewidth]{request-edit}
	\caption{Редактирование заявки}
	\label{request-edit:image}
\end{figure}

\subsubsection{Тестовый сценарий: Изменение статуса заявки}

\textbf{Предусловие:} в системе существуют заявка в статусе "<Черновик">.

\textbf{Тестовый случай:} перевести её в статус "<Назначена">, путем назначения конкретного мастера.

\textbf{Ожидаемый результат:} статус заявки изменяется.

\textbf{Результат:} тест пройден успешно.

\begin{figure}[H]
	\centering
	
\includegraphics[width=0.7\linewidth]{request-status-edit}
	\caption{Изменение статуса заявки}
	\label{request-status-edit:image}
\end{figure}

\subsubsection{Тестовый сценарий: Управление пользователями системы}

\textbf{Предусловие:} пользователь авторизован под учётной записью администратора.

\textbf{Тестовый случай:} администратор открывает раздел "<Пользователи системы"> и просматривает список зарегистрированных сотрудников.

\textbf{Ожидаемый результат:} отображается перечень пользователей с указанием их основных данных и назначенных ролей. Администратор имеет возможность удалять пользователей.

\textbf{Результат:} тест пройден успешно.

\begin{figure}[H]
	\centering
	\includegraphics[width=0.95\linewidth]{users}
	\caption{Интерфейс управления пользователями системы}
	\label{users:image}
\end{figure}

\subsubsection{Тестовый сценарий: Создание нового пользователя}

\textbf{Предусловие:} пользователь авторизован под учётной записью администратора.

\textbf{Тестовый случай:} добавить нового пользователя и назначить ему роль.

\textbf{Ожидаемый результат:} учётная запись появляется в системе.

\textbf{Результат:} тест пройден успешно.

\begin{figure}[H]
	\centering
	\includegraphics[width=0.8\linewidth]{user-create}
	\caption{Интерфейс управления пользователями системы}
	\label{user-create:image}
\end{figure}

\subsubsection{Тестовый сценарий: Просмотр журнала изменений}

\textbf{Предусловие:} в системе были выполнены операции, подлежащие аудиту.

\textbf{Тестовый случай:} пользователь открывает раздел журнала изменений.

\textbf{Ожидаемый результат:} система отображает сведения о выполненных действиях, включая дату и время события, пользователя, тип операции и изменённую сущность.

\textbf{Результат:} тест пройден успешно.

\begin{figure}[H]
	\centering
	\includegraphics[width=0.95\linewidth]{audit}
	\caption{Журнал изменений системы}
	\label{audit:image}
\end{figure}

\subsubsection{Тестовый сценарий: Проверка разграничения прав}

\textbf{Предусловие:} выполнен вход под пользователем с ролью "<Мастер">.

\textbf{Тестовый случай:} попытаться открыть раздел управления пользователями.

\textbf{Ожидаемый результат:} у пользователя отсутствует интерфейская возможность перейти в раздел управления пользователями.

\textbf{Результат:} тест пройден успешно.

\begin{figure}[H]
	\centering
	\includegraphics[width=0.95\linewidth]{master}
	\caption{Главная панель пользователя "<Мастер">}
	\label{master:image}
\end{figure}
